\documentclass[11pt]{article}
\usepackage[review]{acl}
\usepackage{times}
\usepackage{latexsym}
\usepackage{booktabs}
\usepackage{amsmath}
\usepackage{graphicx}
\usepackage{multirow}

\title{Emotion Shifts in Classical Chinese Poetry\\Around the An Lushan Rebellion: A Reproducible Short-Paper Pipeline}

\author{
Anonymous ACL submission
}

\begin{document}
\maketitle

\begin{abstract}
We present a reproducible pipeline for analyzing emotion shifts in Classical Chinese poetry around a major historical event, the An Lushan Rebellion.
Starting from poem-level records, we build sentence-level data, perform emotion annotation, and combine embedding-based clustering with temporal and author-level aggregation.
Our workflow is designed for low-resource literary analysis settings where modern NLP benchmarks do not transfer directly.
We provide a compact short-paper template and an implementation-oriented repository structure to support fast replication and extension.
\end{abstract}

\section{Introduction}

Historical literary corpora offer rich evidence for cultural and emotional dynamics, but they are challenging for contemporary NLP due to archaic language, sparse supervision, and domain-specific semantics.
Classical Chinese poetry is a representative case: emotional expression is dense yet implicit, and large historical events can alter tone, imagery, and rhetorical style.

This paper focuses on emotion shifts around the An Lushan Rebellion and contributes a practical short-paper pipeline that can be reused by beginners and domain researchers.
Compared with a purely narrative digital humanities workflow, our approach emphasizes reproducibility with script-level input/output chains and explicit data transformations.

Our main contributions are:
\begin{itemize}
    \item A sentence-level emotion analysis pipeline for Classical Chinese poetry.
    \item A combined view of emotion through annotation, clustering, and temporal-author analysis.
    \item A submission-ready repository layout for ACL-style short papers.
\end{itemize}

We position this work at the intersection of literary NLP and emotion analysis
\cite{PLACEHOLDER_literary_nlp_survey,PLACEHOLDER_chinese_emotion_resource}.

\section{Task and Data}

\subsection{Problem Formulation}

Given a collection of poems with metadata (author, year, location), our goal is to estimate how emotion distributions change before, during, and after a major historical disruption.
We treat each poem as a set of sentence-level units and analyze both fine-grained emotional labels and unsupervised semantic clusters of emotional statements.

\subsection{Data Processing Pipeline}

Our data pipeline follows three deterministic stages:

\begin{enumerate}
    \item \textbf{Poem-to-sentence conversion}: parse poem text, remove bracketed notes, split by punctuation, and preserve metadata.
    \item \textbf{Emotion annotation}: assign sentence-level emotional tags (with optional multi-rater checks).
    \item \textbf{Aggregation and analysis}: compute year/author/location summaries and run embedding + clustering for latent emotional themes.
\end{enumerate}

The implementation is script-driven and documented with explicit input/output files in the repository.

\section{Method}

\subsection{Sentence-Level Annotation}

We transform each poem into sentence instances and annotate emotion labels for each sentence.
This step yields a structured table suitable for both descriptive statistics and machine-learning analyses.
When two annotators are available, we report agreement indicators and confusion patterns.

\subsection{Embedding and Clustering}

To capture latent emotional semantics beyond manual labels, we encode emotional sentences into dense vectors and cluster them with a fixed number of clusters \(k\).
The cluster outputs are used as interpretable themes and compared across time periods.

\subsection{Temporal and Authorial Analysis}

We aggregate emotion signals by year and author, then inspect changes around the historical event window.
This allows us to separate corpus-level trends from poet-specific expression shifts.

\section{Experiments and Results}

\subsection{Setup}

The current release is built from a reproducible script chain:
\texttt{anshi.py} \(\rightarrow\) \texttt{annotate.py} \(\rightarrow\) \texttt{anshi\_count.py} and
\texttt{emopoem\_get.py} \(\rightarrow\) \texttt{embedding\_get.py} \(\rightarrow\) \texttt{kcluster.py}.

Table~\ref{tab:main_results} is a template for core quantitative outcomes.
Replace placeholder values with final statistics before submission.

\begin{table}[t]
\centering
\small
\begin{tabular}{lcc}
\toprule
Metric & Value & Note \\
\midrule
Sentence instances & \textbf{[TODO]} & from cleaned poems \\
Annotated instances & \textbf{[TODO]} & final labels \\
Inter-annotator agreement & \textbf{[TODO]} & e.g., kappa/F1 \\
Best cluster count \(k\) & \textbf{[TODO]} & elbow + inspection \\
\bottomrule
\end{tabular}
\caption{Main dataset and analysis statistics (to be finalized).}
\label{tab:main_results}
\end{table}

\subsection{Qualitative Findings}

We observe three recurring qualitative patterns (to be backed by final numbers):
\begin{itemize}
    \item Emotion polarity shifts with historical phase transitions.
    \item Some poets show stable style but changing emotional intensity.
    \item Cluster themes reveal rhetorical modes not fully captured by discrete labels.
\end{itemize}

Figure~\ref{fig:cluster_placeholder} is a placeholder for cluster visualization.

\begin{figure}[t]
\centering
\fbox{\rule{0pt}{3.2cm}\rule{0.95\linewidth}{0pt}}
\caption{Placeholder for cluster visualization (replace with real figure).}
\label{fig:cluster_placeholder}
\end{figure}

\section{Discussion, Limitations, and Ethics}

\paragraph{Interpretation Scope.}
Our conclusions depend on available poem metadata and annotation quality.
Historical causality should be interpreted cautiously; we report correlations and temporal alignments rather than definitive causal claims.

\paragraph{Limitations.}
The current version may inherit label bias from prompt design, annotator preference, and domain-specific ambiguity in Classical Chinese syntax.
Future work should evaluate alternative annotation schemes and robust uncertainty estimates.

\paragraph{Ethics and Reproducibility.}
The project uses historical literary texts and focuses on scholarly analysis.
To support reproducibility, we provide script-level pipelines and explicit file dependencies.
Any API-based annotation step should document model version, prompt settings, and cost constraints.

\section{Conclusion}

We provide a practical ACL short-paper blueprint for emotion analysis in Classical Chinese poetry around a major historical event.
The key value is not only the empirical findings but also the reproducible workflow that new researchers can adapt quickly.
The next step is to finalize numeric results, verify references, and submit an anonymized ACL-compliant version.

\section*{Citation Notice}
\small
This draft intentionally uses placeholder citations for safety. Replace all
\texttt{PLACEHOLDER\_*} entries in \texttt{refs.bib} with verified papers before submission.

\bibliography{refs}
\bibliographystyle{acl_natbib}

\end{document}
